\section{A produced analytic contract: rational linear ODE transport}
\label{sec:ode-producer}

\subsection{The precise repository boundary}

The pinned source
\nolinkurl{Derivations/Lean/IPRADerivations/Foundations/OdeRefinement.lean}
proves two arithmetic consequences of an error recurrence: a finite
Gronwall bound and geometric telescoping. Its header explicitly states that
the recurrence describing the Euler polygons is a hypothesis. The corresponding
row of \nolinkurl{Derivations/README.md} also leaves the polygons, their bounds,
interpolation, and the resulting name unconstructed. Thus another abstract
Gronwall interface would not advance this boundary. We supply the missing
producer for nonautonomous linear systems, including rational coefficient
functions and a concrete special-function application.

The result is deliberately quantitative. Every estimate below is an inequality
between finite rational arrays. Its intended object-iPRA formalization uses
signed integer numerators and positive natural denominators, clears the
positive denominators, and applies quantifier-free induction to the resulting
primitive-recursive terms. The source boundary is documented in \cite{repo-ode,repo-gronwall}. The new object derivations are a formalization
milestone; the human proofs below do not assert that those derivations have
already been checked by Lean.

\subsection{Data, norms, and the exact finite program}

Fix a finite dimension $d$. A rational coefficient kernel is a particular
primitive-recursive program evaluating a matrix $B(t)$ at a rational
$t\in[0,1]$. The word ``particular'' matters: we do not assume a total
primitive-recursive evaluator for arbitrary primitive-recursive program codes.
A finite family of such kernels, with a bounded instruction interpreter, is
another admissible implementation.

For real rational vectors use the maximum norm and for matrices the maximum
absolute row-sum norm. For complex rational vectors one may instead take
$\lVert v\rVert=\max_i(|\Re v_i|+|\Im v_i|)$ and the associated row-sum
upper bound. These norms are rational-valued and obey
\[
 \lVert v+w\rVert\le\lVert v\rVert+\lVert w\rVert,\qquad
 \lVert Av\rVert\le\lVert A\rVert\lVert v\rVert,\qquad
 \lVert AB\rVert\le\lVert A\rVert\lVert B\rVert.
\]
For example, expand $(Av)_i=\sum_jA_{ij}v_j$, apply the scalar triangle
inequality, bound each $\lVert v_j\rVert$ by the finite maximum, and take
another finite maximum. The matrix-product inequality follows by applying
the same calculation twice. These proofs are finite sum inequalities.

Assume that natural numbers $K,L$ certify, at all rational arguments in the
interval,
\begin{equation}\label{eq:ode-KL}
 \lVert B(t)\rVert\le K,\qquad
 \lVert B(t)-B(s)\rVert\le L|t-s|.
\end{equation}
Let $v$ be the initial rational vector, and let $U\ge1$ be a natural number
with $\lVert v\rVert\le U$. Put
\begin{equation}\label{eq:ode-constants}
 G=3^K,\qquad Y=GU,\qquad
 H=(L+K^2)Y,\qquad C=GH=(L+K^2)G^2U.
\end{equation}
These constants are intentionally conservative. Their explicit form is more
useful here than a sharper bound involving an unconstructed exponential.

For a positive natural $n$, define the exact Euler array by
\begin{equation}\label{eq:exact-euler}
 X_{n,0}=v,\qquad
 X_{n,j+1}=X_{n,j}+\frac1n B(j/n)X_{n,j}
 \quad(0\le j<n).
\end{equation}
For rational $t<1$, put $j=\lfloor nt\rfloor$ and
\[
 X_n(t)=X_{n,j}+(t-j/n)B(j/n)X_{n,j};
 \qquad X_n(1)=X_{n,n}.
\]
This definition uses rational arithmetic, division with remainder, and a
bounded recurrence. It is a primitive-recursive program uniformly in all
its coded finite data. The integer growth of its intermediate numerators
and denominators does not change this conclusion.

\begin{lemma}[An integer geometric bound]\label{lem:ode-growth}
For positive $n$, natural $K$, and $0\le j\le n$,
\[
 (1+K/n)^j\le3^K.
\]
\end{lemma}
\begin{proof}
First prove, by induction on $j\le n$,
\[
 (1+1/n)^j\le 1+j/n+j^2/n^2.
\]
Multiplying the right side by $1+1/n$ gives
$1+(j+1)/n+(j^2+j)/n^2+j^2/n^3$.
This is at most $1+(j+1)/n+(j+1)^2/n^2$, because
$j^2/n\le j\le j+1$. At $j=n$ the displayed upper bound is $3$.
A separate induction on $K$ proves Bernoulli's inequality
$1+K/n\le(1+1/n)^K$. Raising it to the $n$th power gives
$(1+K/n)^n\le3^K$. Finally $1+K/n\ge1$ and $j\le n$.
Every line can be multiplied by a power of the positive $n$ to become
an integer polynomial inequality. These are the same arithmetic bounds
already available in \nolinkurl{Foundations/DiscreteGronwall.lean}.
\end{proof}

\begin{lemma}[Bounds produced by the Euler program]\label{lem:ode-polybound}
At every rational $s,t\in[0,1]$,
\[
 \lVert X_n(t)\rVert\le Y,\qquad
 \lVert X_n(t)-X_n(s)\rVert\le KY|t-s|.
\]
\end{lemma}
\begin{proof}
Induction in \eqref{eq:exact-euler} gives
$\lVert X_{n,j}\rVert\le(1+K/n)^jU\le GU$.
An interpolated value is a convex combination of its two endpoint values,
so it has the same bound. On the $j$th mesh interval the slope is
$B(j/n)X_{n,j}$ and has norm at most $KY$. Split the rational interval
between $s$ and $t$ at the finitely many mesh points it contains. Apply
this bound to every piece and add their lengths. This proves the second
inequality, including intervals crossing many mesh cells.
\end{proof}

\subsection{Deriving the formerly assumed refinement recurrence}

\begin{theorem}[Complete finite refinement estimate]\label{thm:ode-produced-refinement}
The polygons computed by \eqref{eq:exact-euler} satisfy
\begin{equation}\label{eq:ode-produced-refine}
 \lVert X_n(t)-X_{2n}(t)\rVert\le\frac{C}{2n}
 \quad\text{for every rational }t\in[0,1].
\end{equation}
Consequently, for $s\ge r$,
\begin{equation}\label{eq:ode-produced-cauchy}
 \lVert X_{2^r}(t)-X_{2^s}(t)\rVert
 \le C(2^{-r}-2^{-s})\le C2^{-r}.
\end{equation}
\end{theorem}
\begin{proof}
Write $\delta=1/(2n)$, $t_j=j\delta$,
$V_j=X_n(t_j)$ and $W_j=X_{2n,j}$. Define the exact coarse residual
\[
 r_j=V_{j+1}-V_j-\delta B(t_j)V_j.
\]
If $j=2i$, then $V_j=X_{n,i}$ and the definition of interpolation gives
$r_j=0$. If $j=2i+1$, put $x=X_{n,i}$. The two coarse half-steps both have
increment $\delta B(t_{j-1})x$, while
$V_j=x+\delta B(t_{j-1})x$. Substitution, with no approximation, gives
\begin{equation}\label{eq:ode-residual-identity}
 r_j=\delta\bigl((B(t_{j-1})-B(t_j))x
       -\delta B(t_j)B(t_{j-1})x\bigr).
\end{equation}
By \eqref{eq:ode-KL} and Lemma~\ref{lem:ode-polybound},
\[
 \lVert r_j\rVert\le\delta^2(L+K^2)Y=\delta^2H.
\]
Thus $e_j=\lVert V_j-W_j\rVert$ satisfies the produced recurrence
\begin{equation}\label{eq:ode-produced-recurrence}
 e_0=0,\qquad
 e_{j+1}\le(1+K\delta)e_j+\delta^2H.
\end{equation}
Iterating this inequality and using Lemma~\ref{lem:ode-growth} with $2n$
in place of $n$ yields, for every $j\le2n$,
\[
 e_j\le G\sum_{i<j}\delta^2H
 \le G(2n)\delta^2H=\frac{C}{2n}.
\]
This is exactly the point at which the repository's
\texttt{OdeRefinement.node\_bound} can be instantiated after rational
encoding and clearing denominators. We have now supplied its recurrence
premise, rather than repackaging it.

Both polygons are affine on each interval $[t_j,t_{j+1}]$ of the fine
mesh. Their difference at an intermediate point is a convex combination
of the differences at the two endpoints. The triangle inequality therefore
extends the node estimate to every rational point, proving
\eqref{eq:ode-produced-refine}. Finally add this bound for
$n=2^r,2^{r+1},\ldots,2^{s-1}$. The finite geometric sum is
$\sum_{i=r}^{s-1}2^{-i-1}=2^{-r}-2^{-s}$, proving
\eqref{eq:ode-produced-cauchy}. Its denominator-cleared proof is the
geometric telescoping already available in the repository.
\end{proof}

\subsection{A primitive-recursive name and its exact demands}

For a natural $M\ge1$, let $\ell(M)$ be the least $j\le M$ such that
$M\le2^j$, and put $\ell(0)=0$. This is bounded search. Its search interval
is sufficient because $2^M\ge M$, proved by elementary induction. In
particular, $\ell$ is primitive recursive; it is not an unbounded
``search until the requested error is small.''

\begin{corollary}[A produced regular name]\label{cor:ode-regular-name}
For rational initial data, put $c=\ell(\max(1,C))$ and
\[
 q_p(t)=X_{2^{p+c}}(t).
\]
Then $q$ is primitive recursive, and
$\lVert q_p(t)-q_q(t)\rVert\le2^{-p}$ whenever $q\ge p$.
\end{corollary}
\begin{proof}
Apply \eqref{eq:ode-produced-cauchy} and $C\le2^c$.
\end{proof}

One can also accept a fixed primitive-recursive initial name $(v_r)$ with
\[
 \lVert v_r\rVert\le U,\qquad
 \lVert v_r-v_s\rVert\le2^{-r}+2^{-s}.
\]
We write $X_n(v;t)$ when the initial vector must be displayed.

\begin{lemma}[Initial-data and coefficient stability]\label{lem:ode-stability}
For the same coefficient kernel,
\[
 \lVert X_n(v;t)-X_n(w;t)\rVert\le G\lVert v-w\rVert.
\]
If $B,\widetilde B$ both have norm at most $K$, and
$\lVert B(t)-\widetilde B(t)\rVert\le\eta$ on the interval, then
\[
 \lVert X_n^B(v;t)-X_n^{\widetilde B}(v;t)\rVert
 \le G^2U\eta.
\]
\end{lemma}
\begin{proof}
The initial-data difference obeys the homogeneous Euler recurrence, so the
product bound in Lemma~\ref{lem:ode-growth} applies directly. For the second
claim subtract the recurrences. The two terms are a previous error multiplied
by $I+B(j/n)/n$, and the new term
$(B(j/n)-\widetilde B(j/n))X_n^{\widetilde B}(v;j/n)/n$.
Hence the norm error satisfies
\[
 e_{j+1}\le(1+K/n)e_j+\eta GU/n,\qquad e_0=0.
\]
Iteration gives $e_j\le G(\eta GU)=G^2U\eta$. The common-mesh affine
interpolation argument extends both estimates off the nodes.
\end{proof}

For named initial data and an exact rational coefficient kernel choose
\begin{equation}\label{eq:ode-demands}
 c_v=\ell(4G),\qquad c_m=\ell(4\max(1,C)),\qquad
 q_p(t)=X_{2^{p+c_m}}(v_{p+c_v};t).
\end{equation}
For $q\ge p$, compare first across meshes at the same initial vector and
then across initial vectors at the finer mesh. The estimates give
\[
 \lVert q_p(t)-q_q(t)\rVert
 \le\tfrac14 2^{-p}
      +\tfrac14(2^{-p}+2^{-q})
 \le\tfrac34 2^{-p}.
\]
The backward precision demand on the initial name is therefore the affine
map $p\mapsto p+c_v$. The mesh resource is the explicit number
$2^{p+c_m}$. This fits the existing numerical precision/cost machinery.

If coefficients themselves arrive through names, use fixed rational
approximating kernels $B_r$, with the same certified $K,L$ and a proved
uniform bound
$\lVert B_r-B_s\rVert\le2^{-r}+2^{-s}$. Allocate all three sources of
error by replacing the $4$ in the mesh and initial shifts with $8$, and
use the coefficient shift $c_B=\ell(8G^2U)$. Mesh error then contributes
at most $2^{-p}/8$, and each of the two input-name comparisons at most
$2^{-p}/4$. Their sum is below $2^{-p}$. The existence of such $B_r$,
including uniform denominator guards, must be proved from the particular
coefficient representation; the stability lemma does not invent it.

\subsection{A path-clearance producer for the Gauss equation}

We now construct the bounds \eqref{eq:ode-KL} from concrete data. This is
useful for the path-continuation theorem in
\texttt{analysis/hypergeometric\_ipra\_2026-11-16/sections/10\_continuation.tex}.
That packet supplies the intended special-function interface; our producer
provides a fully explicit finite arithmetic alternative to its local
Taylor-majorant step.

Let $z(t)=u+th$ be one complex-rational line segment, and consider a singular
point $s\in\{0,1\}$. Put $a=u-s$, $E=|h|_2^2$, and
$D=\Re(\overline a h)$, where $|x+iy|_2^2=x^2+y^2$ involves no square
root. If $E=0$, set $\mu_s=|a|_2^2$. Otherwise set
\[
 \tau=\min(1,\max(0,-D/E)),\qquad
 \mu_s=|a+\tau h|_2^2.
\]

\begin{lemma}[Exact segment clearance]\label{lem:ode-clearance}
The rational number $\mu_s$ is the minimum squared distance from the
segment to $s$. The test $\mu_s>0$ succeeds exactly when the segment misses
$s$, and supplies a lower bound valid at every rational parameter $t$.
\end{lemma}
\begin{proof}
The distance polynomial is $f(t)=|a|_2^2+2Dt+Et^2$. If $E=0$ it is constant.
If $-D/E\in[0,1]$, completing the square gives
$f(t)-f(-D/E)=E(t+D/E)^2\ge0$. If $-D/E<0$, then $D>0$ and
$f(t)-f(0)=t(2D+Et)\ge0$. If $-D/E>1$, then $D<-E$, and for $t\le1$,
$f(t)-f(1)=(t-1)(2D+E(t+1))\ge0$. The boundary cases satisfy either
adjacent formula. Thus the minimum is achieved at the displayed rational
$\tau$. It is zero precisely when $a+\tau h=0$.
\end{proof}

Assume both clearance tests pass. Write $|z|_1=|\Re z|+|\Im z|$ and set
\[
 R=\max(|u|_1,|u+h|_1),\qquad T=|h|_1,\qquad
 D_*=\frac{R(1+R)}{\mu_0\mu_1}.
\]
Convexity gives $|z(t)|_1\le R$. Since
$|1/z|_1=|z|_1/|z|_2^2$, the clearances imply
\begin{equation}\label{eq:ode-inverse-bound}
 \left|\frac1{z(1-z)}\right|_1\le D_*.
\end{equation}
For two points $z,w$ on the segment, putting $d(z)=z(1-z)$ gives
\[
 d(z)-d(w)=(z-w)(1-z-w),\qquad
 |d(z)-d(w)|_1\le(1+2R)|z-w|_1.
\]
The reciprocal difference identity therefore gives
\begin{equation}\label{eq:ode-inverse-lip}
 \left|\frac1{d(z)}-\frac1{d(w)}\right|_1
 \le D_*^2(1+2R)|z-w|_1.
\end{equation}
All tests, constants, and proofs so far are rational arithmetic.

For complex-rational hypergeometric parameters $a,b,c$, put
\[
 S=|a|_1+|b|_1+1,\quad Q=|a|_1|b|_1,\quad
 P=SR+|c|_1,\quad N=Q+P.
\]
The Gauss equation \cite{dlmf-gauss} becomes a first-order system in $Y=(y,y')$:
\[
 \frac{dY}{dz}=A(z)Y,\qquad
 A(z)=\begin{pmatrix}
 0&1\\[1mm]
 \displaystyle\frac{ab}{d(z)}&
 \displaystyle\frac{(a+b+1)z-c}{d(z)}
 \end{pmatrix}.
\]
On the segment the coefficient matrix is $B(t)=hA(z(t))$. Define natural
numbers by exact rational ceilings:
\begin{align}
 K&=\left\lceil T(1+ND_*)\right\rceil,\label{eq:gauss-K}\\
 L&=\left\lceil T^2\bigl(SD_*+ND_*^2(1+2R)\bigr)\right\rceil.
 \label{eq:gauss-L}
\end{align}

\begin{proposition}[Produced Gauss transport data]\label{prop:gauss-transport-data}
The matrices $B(t)$ satisfy \eqref{eq:ode-KL} with
\eqref{eq:gauss-K}--\eqref{eq:gauss-L}. Consequently a certified rational
segment avoiding $0,1$, together with initial values, produces the regular
name and explicit precision schedule of \eqref{eq:ode-demands}.
\end{proposition}
\begin{proof}
The first row of $B$ has row norm $T$. The second has row norm at most
$TND_*$ by \eqref{eq:ode-inverse-bound}; their maximum is bounded by
$T(1+ND_*)$. For differences, the first row is constant. For the second
row use \eqref{eq:ode-inverse-lip}; the two constant numerator bounds sum
to $N$, and the varying numerator has Lipschitz bound $S$ in $z$.
Thus the row difference is at most
$T\bigl(SD_*+ND_*^2(1+2R)\bigr)|z-w|_1$.
Since $|z(t)-z(s)|_1=T|t-s|$, \eqref{eq:gauss-L} follows.
\end{proof}

\begin{proposition}[Producing coefficient-name stability from parameter names]
\label{prop:gauss-parameter-stability}
Suppose two parameter triples $(a,b,c)$ and $(\widetilde a,\widetilde b,
\widetilde c)$ have all six $|\cdot|_1$ norms at most a natural $A$, and
corresponding entries differ by at most $\varepsilon$. For the same segment,
\[
 \lVert B(t)-\widetilde B(t)\rVert\le J\varepsilon,
 \qquad
 J=\left\lceil TD_*(2A+2R+1)\right\rceil.
\]
The bounds \eqref{eq:gauss-K}--\eqref{eq:gauss-L} hold uniformly for every
such triple after using
$S=2A+1$, $Q=A^2$, $P=(2A+1)R+A$, and $N=Q+P$.
\end{proposition}
\begin{proof}
The denominator depends only on the fixed path. The identity
$ab-\widetilde a\widetilde b=a(b-\widetilde b)+
\widetilde b(a-\widetilde a)$ bounds the first numerator difference by
$2A\varepsilon$. The second numerator difference is
$((a-\widetilde a)+(b-\widetilde b))z-(c-\widetilde c)$ and has norm at most
$(2R+1)\varepsilon$. Add the two row contributions and multiply by $TD_*$.
The uniform majorants follow directly from the common parameter bound.
\end{proof}

This proposition supplies an actual coefficient-name producer. If $a_r,b_r,c_r$
are three particular PR regular names with all approximants bounded by $A$,
substitution in the rational Gauss matrix gives PR kernels with the proved
common $K,L$ and the proved comparison bound $J(2^{-r}+2^{-s})$.
For a single segment, a safe joint schedule for the parameter names, the
initial name, and the Euler mesh uses
\[
 c_{\mathrm{par}}=\ell(8G^2U\max(1,J)),\qquad
 c_v=\ell(8G),\qquad c_m=\ell(8\max(1,C)).
\]
At output precision $p$, query all three parameters at $p+c_{\mathrm{par}}$,
the initial value at $p+c_v$, and use $2^{p+c_m}$ cells.
Comparing two outputs uses, in order, the mesh error, the initial-data error,
and the parameter error. Their upper bounds for $q\ge p$ are respectively
$2^{-p}/8$, $(2^{-p}+2^{-q})/8$, and
$(2^{-p}+2^{-q})/8$, whose sum is at most $5\cdot2^{-p}/8$.
Thus parameter stability is not left as an uninhabited analytic premise.

This ordinary-point transport works for resonant parameters too. Producing
the particular initial values of ${}_2F_1(a,b;c;z)$ from its regular series
still requires the series packet's parameter and denominator conditions.
The transport does not establish gamma connection coefficients, contour
bounds, or modular inversion. Those are different producer obligations.

\subsection{A whole finite path with one closed bound}

Let a finite rational polygonal path have $m$ certified edges, with data
$K_j,L_j$ and $G_j=3^{K_j}$. For $m=0$ the output is the input. For $m>0$
put
\begin{equation}\label{eq:ode-path-constants}
 G_{\mathrm{tot}}=\prod_{j<m}G_j,\qquad
 C_{\mathrm{path}}=
 U G_{\mathrm{tot}}\sum_{j<m}G_j(L_j+K_j^2).
\end{equation}
Apply the Euler program with the same $n=2^r$ subdivisions on each edge,
feeding the resulting rational endpoint into the next edge.

\begin{theorem}[Finite-path endpoint name]\label{thm:ode-path-name}
For rational initial data, consecutive global refinements have discrepancy
at most $C_{\mathrm{path}}/(2\cdot2^r)$, and the tail from level $r$ to any
finer level is at most $C_{\mathrm{path}}2^{-r}$. Thus a finite path has a
primitive-recursive endpoint name. For named initial data, valid shifts are
$\ell(4G_{\mathrm{tot}})$ for the input and
$\ell(4\max(1,C_{\mathrm{path}}))$ for the mesh exponent.
\end{theorem}
\begin{proof}
Let $E_{j,n}$ be the linear endpoint map of the $j$th Euler program.
Its operator norm is at most $G_j$. The proof of
Theorem~\ref{thm:ode-produced-refinement}, applied to the matrix initial value $I$ with row-sum norm at most one, gives
\[
 \lVert E_{j,n}-E_{j,2n}\rVert
 \le\frac{G_j^2(L_j+K_j^2)}{2n}.
\]
The difference of two finite products is the sum obtained by changing one
factor at a time. The $j$th term, acting on $v$, has norm at most
\[
 U\left(\prod_{\ell>j}G_\ell\right)
 \frac{G_j^2(L_j+K_j^2)}{2n}
 \left(\prod_{\ell<j}G_\ell\right)
 =\frac{UG_{\mathrm{tot}}G_j(L_j+K_j^2)}{2n}.
\]
Add these finitely many bounds and then telescope dyadic levels. The product
also has initial-data sensitivity $G_{\mathrm{tot}}$, giving the two named
input shifts by the proof of \eqref{eq:ode-demands}.
\end{proof}

For named Gauss parameters shared along the entire path, let $J_j$ be the
produced parameter Lipschitz constant on edge $j$, and put
\[
 S_{\mathrm{path}}=UG_{\mathrm{tot}}\sum_{j<m}G_jJ_j.
\]
The same one-factor-at-a-time product comparison, now using
Lemma~\ref{lem:ode-stability} rather than the mesh estimate, bounds the
endpoint change by $S_{\mathrm{path}}\varepsilon$ when the three parameter
changes are at most $\varepsilon$. A complete joint path schedule therefore
uses shifts $\ell(8\max(1,S_{\mathrm{path}}))$ for the parameter names,
$\ell(8G_{\mathrm{tot}})$ for the initial name, and
$\ell(8\max(1,C_{\mathrm{path}}))$ for the common mesh exponent. All constants
are finite products and sums over the supplied path.

\subsection{The same producer for higher-order equations and Meijer \texorpdfstring{$G$}{G}}

The construction is not restricted to two-dimensional systems. A scalar equation
\[
 a_q(z)y^{(q)}+\sum_{j<q}a_j(z)y^{(j)}=0,\qquad q\ge1,
\]
with polynomial coefficients becomes a rational companion system for
$(y,y',\ldots,y^{(q-1)})$ whenever a finite certificate bounds the reciprocal
of $a_q$ on the path. The superdiagonal entries are one and the last row is
$-a_j/a_q$. Polynomial coefficient sums and the finite power-difference
identity compute bounds for each numerator and its Lipschitz constant.
The reciprocal-difference identity supplies the remaining denominator
estimate. Thus this is one fixed PR kernel with numerical order and
coefficient lists, not a separate implementation for every equation.

For a concrete class, let $\theta=z\,d/dz$ and
$\kappa=(-1)^{p-m-n}\in\{-1,1\}$. The Meijer $G$ equation \cite{dlmf-meijer} is
\[
 \left(\kappa z\prod_{i=1}^{p}(\theta-a_i+1)
       -\prod_{j=1}^{q}(\theta-b_j)\right)y=0
.
\]
Here $p,q,m,n$ are natural numbers with $0\le m\le q$, $0\le n\le p$; take $p\le q$ and $q\ge1$. Expand the two finite products as
polynomials in $\theta$. The identity
\[
 \theta^k=\sum_{j=0}^k S(k,j)z^j\frac{d^j}{dz^j},\qquad
 S(k+1,j)=jS(k,j)+S(k,j-1),
\]
with $S(0,0)=1$ and out-of-range entries zero, converts them to ordinary
derivatives. To prove it, apply $z\,d/dz$ to each summand and collect the
two resulting coefficients. This finite recurrence is PR and its coefficient
identity has a QF-induction proof. The highest coefficient is
\[
 a_q(z)=
 \begin{cases}-z^q,&p<q,\\z^q(\kappa z-1),&p=q.\end{cases}
\]
For a rational segment the exact clearance procedure above therefore needs
only the point $0$, and additionally $\kappa$ when $p=q$.

Here are explicit companion bounds. Let $R$ bound $|z|_1$ on the segment,
$T=|h|_1$, and $\mu_s>0$ be its minimum squared distance to each needed
singularity. A reciprocal bound is
\[
 D_q=(R/\mu_0)^q\quad(p<q),\qquad
 D_q=(R/\mu_0)^q(R+1)/\mu_\kappa\quad(p=q).
\]
Write $a_j(z)=\sum_k c_{jk}z^k$, and compute
$P_j=\sum_k|c_{jk}|_1R^k$ and
$L_j=\sum_{k\ge1}k|c_{jk}|_1R^{k-1}$, including $j=q$.
Then $|a_j(z)|_1\le P_j$, its difference is bounded by
$L_j|z-w|_1$, and the reciprocal difference is bounded by
$D_q^2 L_q|z-w|_1$. Consequently valid natural constants are
\[
 K=\left\lceil T\left(1+D_q\sum_{j<q}P_j\right)\right\rceil,
\]
\[
 L=\left\lceil T^2\left(D_q\sum_{j<q}L_j+
                 D_q^2 L_q\sum_{j<q}P_j\right)\right\rceil.
\]
These formulas, the finite expansion, and the clearance tests produce all
the rational transport data from finite rational parameters. Repeated
parameters do not invalidate an ordinary-point companion equation.
They can, however, change the difficulty of constructing the initial jet
of the particular Meijer $G$ branch. This section supplies continuation
\emph{from a supplied certified jet}; it does not identify that jet with a
Mellin--Barnes contour integral or prove a connection formula.

\subsection{Integral observations, uniqueness, and the semantic boundary}

A Cauchy estimate alone should not be advertised as an ODE theorem. Here the
finite output also satisfies the required integral observations. Let
$F_n(t)=B(t)X_n(t)$. From \eqref{eq:ode-KL} and
Lemma~\ref{lem:ode-polybound},
\[
 \lVert F_n(t)-F_n(s)\rVert\le H|t-s|.
\]
On the $j$th Euler cell the polygon slope is the constant
$Q_n(t)=B(j/n)X_{n,j}$, and
\begin{equation}\label{eq:ode-local-residual}
 \lVert F_n(t)-Q_n(t)\rVert\le H/n.
\end{equation}
The finite integral of $Q_n$ is exactly $X_n(t)-v$, by the defining recurrence
and the last partial cell. For an independent positive mesh number $M$, define
\[
 R_M(F_n;t)=\frac{t}{M}\sum_{i<M}F_n(it/M).
\]
A common rational refinement of the Euler cells and the $M$ quadrature cells
proves
\begin{equation}\label{eq:ode-integral-observation}
 \left\lVert X_n(t)-v-R_M(F_n;t)\right\rVert
 \le H/n+H/M.
\end{equation}
Indeed, replace $F_n$ on each quadrature cell by its left value. Its variation
is at most $H/M$, and the sum of all cell lengths is at most one. Next use
\eqref{eq:ode-local-residual} on the common refinement to replace $F_n$ by
$Q_n$; that contributes at most $H/n$. This is a proof about two finite sums,
not an invocation of integration or completeness inside iPRA. Choosing
$n,M$ to make the two rational upper bounds at most a prescribed dyadic
error produces the integral-equation observation contract.

The uniform $KY$-Lipschitz bound on all Euler polygons passes to their finite observation laws. Thus the pointwise names are spatially coherent and represent one uniformly continuous function. In a semantic interpretation of regular names, the same estimates identify
the output with the solution of
\[
 X(t)=v+\int_0^tB(s)X(s)\,ds.
\]
For completeness, uniqueness has a constructive quantitative proof. Suppose
a second solution $Z$ has the same initial value and a bound
$\lVert Z(t)\rVert\le V$. Its integral equation implies that it is
$KV$-Lipschitz. Hence $B(t)Z(t)$ is $(L+K^2)V$-Lipschitz. On a mesh cell,
subtract its left rectangular contribution from its integral; the remaining
norm is at most $(L+K^2)V/(2n^2)$. Comparison with Euler and finite Gronwall
therefore gives
\[
 \lVert Z(j/n)-X_{n,j}\rVert
 \le\frac{G(L+K^2)V}{2n}.
\]
The interpolation estimate gives convergence off the nodes as well. For
every requested precision choose the displayed finite $n$ and the necessary
finite integral observations; the two solution names must agree at that
precision. No zero decision or infinite subsequence is selected.

Splitting an edge at a rational fraction $\theta$ replaces $B(t)$ on the
first subinterval by $\theta B(\theta t)$ and on the second by
$(1-\theta)B(\theta+(1-\theta)t)$. The corresponding bounds scale as
$K\mapsto\theta K$, $L\mapsto\theta^2L$ before rounding up. Adding the two
finite integral observations recovers the original one. The uniqueness
estimate proves equality of the endpoint names under rational subdivision.
This does not imply equality for different paths around singularities;
monodromy is compatible with every theorem above.

The formal object statement quantifies natural codes and precision variables
for particular compiled kernels. It does not quantify all real functions.
The semantic statement that such names denote classical or constructive
solutions is a separate interpretation theorem. Keeping this boundary visible
both respects iPRA's language and preserves the substantive analytic content.

\subsection{Lean implementation plan at the actual source boundary}

The following dependency order produces concrete endpoints instead of a
new layer of uninhabited contracts.

\begin{enumerate}
\item Encode signed rational arrays, positive denominators, finite maximum
and row-sum norms, and guarded inverse. Establish the displayed finite norm
inequalities as object-theory lemmas.
\item Define the primitive-recursive Euler node and interpolation programs.
Derive their equations, totality of the codes, and denominator positivity.
\item Prove the node growth and polygon Lipschitz bounds. Reuse the existing
\nolinkurl{DiscreteGronwall.growth}; induction invariants are denominator-cleared
quantifier-free inequalities.
\item Prove \eqref{eq:ode-residual-identity} by splitting the fine-node index
into its parity cases. Bound its two terms to obtain
\eqref{eq:ode-produced-recurrence}. This is the central new producer lemma.
\item Instantiate \texttt{OdeRefinement.node\_bound} and
\texttt{geometric\_cauchy}, then prove the affine interpolation estimate and
the explicit $\ell$-based name schedule.
\item Prove initial-data and coefficient stability. Expose the exact affine
precision demands to the already existing precision/cost graph interface.
\item Implement segment clearance by the three rational clamping cases and
bind the Gauss matrix to \eqref{eq:gauss-K}--\eqref{eq:gauss-L}.
\item Prove finite product telescoping and obtain
\eqref{eq:ode-path-constants} uniformly in coded path length. Connect the
initial-value producer supplied by the particular special-function packet.
\item Formalize the common-refinement sum proof of
\eqref{eq:ode-integral-observation}. Treat the semantic solution
interpretation and the object-level finite observation theorem as separately
named declarations.
\end{enumerate}

For a fixed PR coefficient kernel, the decisive eventual theorem has literal
judgment \texttt{Prf PRAAx}: from denominator guards and the rational
coefficient bounds it derives the coded Euler refinement inequality. A
second literal \texttt{Prf PRAAx} theorem assembles the named outputs. A
host Lean theorem about \texttt{Real}, a numerical run, or successful
elaboration of an interface structure does not substitute for either
endpoint.
